\documentclass[conference]{IEEEtran}
\usepackage{times}
\usepackage{cite}
\usepackage{amsmath,amssymb,amsfonts}
\usepackage{graphicx}
\usepackage{booktabs}
\usepackage{url}
\usepackage{float}
\setlength{\textfloatsep}{6pt}
\setlength{\intextsep}{6pt}
\setlength{\abovecaptionskip}{4pt}

% exact page geometry per specification (US Letter)
\setlength{\paperwidth}{8.5in}
\setlength{\paperheight}{11in}
\setlength{\textwidth}{7.0in}        % 2 x 3.5in columns
\setlength{\textheight}{9.25in}       % 11 - 0.75 top - 1.0 bottom
\setlength{\columnsep}{0.25in}
\setlength{\oddsidemargin}{-0.375in}  % 1in - 0.375in = 0.625in left
\setlength{\evensidemargin}{-0.375in}
\setlength{\topmargin}{-0.25in}       % 1in - 0.25in = 0.75in top
\setlength{\headheight}{0pt}
\setlength{\headsep}{0pt}
\setlength{\footskip}{0.5in}
\setlength{\parindent}{0.2in}
\setlength{\parskip}{0pt}
\pagestyle{empty}

% balance the two columns on the final page
\IEEEtriggeratref{9}
\begin{document}

\title{\fontsize{24}{28}\selectfont The Open-Weight Oligopoly: Measuring
Concentration and Adoption in the Hugging Face Model Ecosystem}

\author{\IEEEauthorblockN{\fontsize{11}{13}\selectfont Saurav Rijal}
\IEEEauthorblockA{\fontsize{10}{12}\selectfont\textit{Department of Computer
Science, Texas State University}\\ San Marcos, TX, USA\\
sauravrijal@txstate.edu}
}
\maketitle
\thispagestyle{empty}

\begin{abstract}
We audit concentration and adoption in the Hugging Face model ecosystem from a reproducible snapshot of the 20,000 most-downloaded models on the public Hub API, with full metadata for the top 2,000. Downloads are extraordinarily concentrated across models (Gini $0.92$) but not across publishers (4,144 organizations; HHI $375$): a handful of artifacts absorb most usage while thousands of publishers compete. Controlling for publisher, license, and task, larger models are downloaded slightly \emph{less}, while corporate publishers and research labs substantially outperform community accounts and restrictive licenses underperform permissive ones. Breakout models are modestly predictable from metadata alone (AUC $0.72$), with model age the dominant feature, and $k$-means separates breakout flagships from an aging mid-tier and two long tails. Cohort analysis shows the top-ten publishers' download share falling from $79\%$ (2022) to $54\%$ (2026) even as median model size rose roughly a hundredfold. The open-weight ecosystem is thus a model-level oligopoly inside an organizationally competitive market: open weights do not imply decentralized influence.
\end{abstract}

\begin{IEEEkeywords}
open-weight models, Hugging Face, market concentration, technology adoption,
machine learning, open source AI
\end{IEEEkeywords}

\section{Introduction}
The Hugging Face Hub has become the default public registry for open-weight
artificial intelligence models, hosting millions of model repositories that
underpin research, startups, and production systems worldwide. ``Open
weights'' now sit at the center of AI policy debates, capability-diffusion
arguments, and the business models of the largest AI labs. Yet the openness
of the \emph{artifact} (weights anyone may download) should not be confused
with the openness of the \emph{ecosystem}: who publishes the models everyone
actually uses, and what separates the handful of breakout models from the
long tail nobody downloads?

This paper audits the Hub's model ecosystem with its own public data. We
collect the 20,000 most-downloaded models via the Hub's REST API on September
29, 2026, and fetch full metadata, including parameter counts, for the top
2,000. We ask: (i) how concentrated are downloads across models and
organizations? (ii) does the ecosystem centralize or diversify over time?
(iii) what attributes (organization type, license class, age, task, model
size) predict adoption, and can breakout models be predicted from metadata
alone? (iv) do models fall into natural archetypes, and where does licensing
fit in?

Our contributions are a reproducible ingest--clean--model pipeline over the
public Hub API with assertion-gated collection; concentration measurement
with the Gini coefficient and the Herfindahl--Hirschman Index (HHI); an
ordinary least squares (OLS) attribution of downloads with robust inference;
a head-to-head comparison of logistic regression and random-forest breakout
classifiers under repeated cross-validation; a validated $k$-means
segmentation of model archetypes; and a cohort analysis of centralization and
licensing since 2022. The headline finding is that the open-weight ecosystem
is, by downloads, an oligopoly: a tiny fraction of models and organizations
captures the overwhelming majority of usage, and breakout adoption is
predictable from a few structural features.

\section{Related Work}
The most directly comparable study is the recent ``Economies of Open
Intelligence'' analysis, which reconstructs the complete history of weekly
Hub downloads from June 2020 to August 2025 across 851,000 models and
documents a rebalancing of power from US industry labs toward unaffiliated
developers, community organizations, and Chinese industry \cite{econ-open-intel}.
That work studies the full population over time; we study a current snapshot
of the top 20,000 models by trailing-30-day downloads and add inferential and
predictive modeling of what \emph{explains} and \emph{predicts} adoption at the
top of the distribution, which the historical census does not attempt.

A systematic literature review with quantitative validation examines claims
about Hugging Face popularity, defining popularity by downloads and studying
the top 1,000 models across snapshots \cite{hf-lit-review}. Our sample is
twenty times larger, uses current data, and adds concentration indices plus
machine-learning classifiers rather than claim validation alone. Licensing on
the Hub has been studied empirically with attention to documentation and
supply-chain challenges, including a manual audit of ``Other''-licensed
models \cite{hf-licensing}; we instead treat license class as a structured
covariate in adoption models and track its composition over time. Recent work
on public application transformation through Model--Space links finds that
downstream reuse is highly concentrated among few models and that downloads,
likes, and reuse capture distinct facets of impact \cite{hf-reuse}, a finding
our download-concentration results corroborate from the artifact side. None of
these joins concentration measurement with predictive modeling of breakout
adoption on a current top-of-Hub sample, which is the gap this paper fills.

\section{Methodology}
\subsection{Data Sources}
Model listings come from the public Hugging Face Hub REST API
(\texttt{huggingface.co/api/models}), queried without authentication on
September 29, 2026 \cite{hf-api}. We page through \texttt{limit=100},
\texttt{sort=downloads}, \texttt{direction=-1} with cursor pagination for 200
pages, yielding the 20,000 models with the highest trailing-30-day download
counts. The \texttt{downloads} field is the rolling 30-day count, distinct
from the cumulative \texttt{downloadsAllTime} available only on request
\cite{hf-hub-docs}; every ``download'' below means trailing-30-day downloads
unless stated otherwise. For the top 2,000 models by downloads we fetch the
per-model detail endpoint, which adds the \texttt{safetensors.total}
parameter count, the author field, and gating status.

\subsection{Cleaning and Feature Engineering}
The publisher is parsed as the \texttt{org/model} id prefix (legacy
single-component ids keep the id as their own org). Organization type is
author-coded from public knowledge into \emph{corporate} (for-profit
companies and their labs), \emph{research} (academic and nonprofit labs), and
\emph{community} (individuals and community converter/quantization groups);
uncoded orgs default to community, and the full mapping is fixed in the
pipeline. License class is rule-parsed from \texttt{license:*} tags, first
match in order: \emph{restrictive} (behavioral-use: \texttt{openrail*},
\texttt{llama*}, \texttt{gemma}, non-commercial, \texttt{apple-amlr},
\texttt{openmdw}), \emph{copyleft} (\texttt{gpl*}, \texttt{agpl*},
\texttt{cc-by-sa*}), \emph{permissive} (\texttt{apache-2.0}, \texttt{mit},
\texttt{bsd-*}, \texttt{cc-by-*}, \texttt{cc0}, \texttt{wtfpl},
\texttt{unlicense}), \emph{other} (any remaining license tag), and
\emph{unspecified} (no license tag). Model age is measured in days from
\texttt{createdAt} to collection; task comes from \texttt{pipeline\_tag}
(69.6\% coverage); parameter counts cover the detail subsample and are
reported with coverage everywhere they are used.

\subsection{Models}
M1 measures download concentration: the Gini coefficient over models and the
HHI over organization download shares,
\begin{equation}
\mathrm{HHI} = 10000 \sum_j s_j^2,
\end{equation}
where $s_j$ is org $j$'s share of sample downloads, alongside CR1/CR4/CR10.
M2 is an OLS regression of $\log(\mathrm{downloads})$ on organization type,
license class, $\log(\mathrm{age})$, and task indicators (top tasks plus
other), with HC3 heteroskedasticity-robust standard errors; a second
specification on the detail subsample adds $\log(\mathrm{parameters})$.
M3 is a logistic regression for $P(\mathrm{breakout})$, where breakout means
downloads at or above the sample 90th percentile, reporting odds ratios with
95\% confidence intervals and $p$-values. M4 is a random-forest breakout
classifier (400 trees) compared honestly against the logistic model by
repeated stratified 5-fold cross-validation (10 repeats, 50 folds), reporting
mean $\pm$ SD area under the curve (AUC) against a majority-class baseline.
M5 is $k$-means on standardized $[\log \mathrm{downloads}, \log \mathrm{likes},
\mathrm{age}, \log \mathrm{org\ size}]$ with silhouette validation over
$k=2,\dots,6$. Every pipeline script is assertion-gated and exits nonzero on
any data mismatch.

\section{Results}
\subsection{Download concentration}
The 20,000 models accumulated 2,835,223,827 trailing-30-day downloads
(median 6,719, mean 141,761, maximum 242,757,705 for
sentence-transformers/all-MiniLM-L6-v2). The download Gini coefficient is 0.9182, and the
top 100 models capture 48.4\% of sample downloads while the top
1,000 capture 81.2\%; Fig.~\ref{fig:lorenz} shows the Lorenz
curve. Table~\ref{tab:descriptives} summarizes the sample.

\begin{figure}[t]
\centering
\includegraphics[width=\columnwidth]{figures/fig1_lorenz.png}
\caption{Lorenz curve of trailing-30-day downloads across the 20,000 sampled
models (Gini 0.9182). The dashed line is perfect equality.}
\label{fig:lorenz}
\end{figure}

\begin{table}[t]
\caption{Download distribution across the 20,000-model sample}
\label{tab:descriptives}
\centering
\small
\begin{tabular}{lr}
\toprule
Statistic & Value \\
\midrule
Models & 20000 \\
Organizations & 4144 \\
Total 30-day downloads & 2,835,223,827 \\
Median downloads & 6,718 \\
Interquartile range & 3,262--27,479 \\
Mean downloads & 141,761 \\
Maximum downloads & 242,757,705 \\
Gini coefficient & 0.9182 \\
Share of top 100 models & 48.4\% \\
Share of top 1,000 models & 81.2\% \\
\bottomrule
\end{tabular}
\normalsize
\end{table}


Aggregating downloads by publisher, 4,144 organizations appear in
the sample. The largest, sentence-transformers, holds 12.0\% of downloads;
the top 4 hold 32.7\% and the top 10 47.7\%
(Table~\ref{tab:toporgs}, Fig.~\ref{fig:toporgs}). The HHI is 375,
which sits below the 1,500 ``moderately concentrated'' threshold
used in merger review \cite{doj-hhi}: by downloads, no single publisher
dominates, but the model-level distribution is far more concentrated than the
organization-level one.

\begin{figure}[t]
\centering
\includegraphics[width=\columnwidth]{figures/fig2_orgs.png}
\caption{Top ten organizations by share of sample downloads (trailing 30
days).}
\label{fig:toporgs}
\end{figure}

\begin{table}[t]
\caption{Top ten organizations by 30-day downloads}
\label{tab:toporgs}
\centering
\small
\begin{tabular}{lrrr}
\toprule
Organization & Models & Downloads & Share \\
\midrule
sentence-transformers & 66 & 339,894,578 & 12.0\% \\
Qwen & 312 & 296,011,816 & 10.4\% \\
BAAI & 34 & 152,501,848 & 5.4\% \\
google & 248 & 140,026,994 & 4.9\% \\
cross-encoder & 31 & 100,287,358 & 3.5\% \\
Comfy-Org & 87 & 82,917,785 & 2.9\% \\
unsloth & 594 & 65,252,240 & 2.3\% \\
openai & 24 & 63,941,670 & 2.3\% \\
facebook & 305 & 59,771,633 & 2.1\% \\
google-bert & 13 & 52,459,563 & 1.8\% \\
\bottomrule
\end{tabular}
\normalsize
\end{table}

\begin{table}[t]
\caption{Concentration of 30-day downloads across organizations}
\label{tab:hhi}
\centering
\small
\begin{tabular}{lr}
\toprule
Measure & Value \\
\midrule
Organizations & 4144 \\
HHI (download shares) & 375 \\
Largest org share (CR1) & 12.0\% \\
Top-4 share (CR4) & 32.8\% \\
Top-10 share (CR10) & 47.7\% \\
\bottomrule
\end{tabular}
\normalsize
\end{table}


\subsection{Parameter scale}
Parameter counts are available for 1,237 of the top 2,000
models (61.9\% of the detail subsample, 6.2\% of the full frame). The median model has 1.02~billion parameters; the parameter Gini
is 0.9021 and the ten largest models hold 29.6\% of
subsample parameters. Median size rises from 124~million in 2022 to 12.0~billion in 2026 (Fig.~\ref{fig:temporal}).

\subsection{What predicts downloads}
Table~\ref{tab:ols} reports the full-sample OLS for $\log(\mathrm{downloads})$
($R^2=0.148$, $n=20{,}000$). Relative to community publishers,
corporate orgs average $e^{1.26}\approx 3.5\times$ the downloads of community orgs and research orgs $e^{0.88}\approx 2.4\times$ as many (both $p<0.001$); relative to permissive licenses,
restrictive (behavioral-use) licenses sit 31\% lower ($\hat\beta=-0.374$, $p<0.001$) and unspecified licenses 8\% lower. Older models accumulate more downloads
(elasticity $0.068$, $p<0.001$). On the detail subsample,
$\log(\mathrm{parameters})$ enters with elasticity $-0.025$ ($p=0.043$), and $R^2$ falls to $0.126$ (Table~\ref{tab:olsparams}).

\begin{table}[t]
\caption{OLS: $\log(\text{downloads})$ on model attributes (HC3 robust SE; key coefficients)}
\label{tab:ols}
\centering
\small
\begin{tabular}{lrrr}
\toprule
Variable & $\hat\beta$ & SE & $p$-value \\
\midrule
const & 8.4327 & 0.0491 & 0.0000 \\
org\_type\_corporate & 1.2593 & 0.0465 & 0.0000 \\
org\_type\_research & 0.8842 & 0.1096 & 0.0000 \\
license\_class\_restrictive & -0.3740 & 0.0339 & 0.0000 \\
license\_class\_unspecified & -0.0787 & 0.0294 & 0.0073 \\
tag\_group\_image-text-to-text & 0.9203 & 0.0444 & 0.0000 \\
tag\_group\_automatic-speech-recognition & 1.3295 & 0.0901 & 0.0000 \\
log\_age & 0.0681 & 0.0082 & 0.0000 \\
\bottomrule
\multicolumn{4}{l}{$n=20000$; $R^2=0.148$; adj.\ $R^2=0.147$. Task dummies n.s.\ omitted.} \\
\end{tabular}
\normalsize
\end{table}

\begin{table}[t]
\caption{OLS with $\log(\text{parameters})$: detail subsample (HC3 robust SE; key coefficients)}
\label{tab:olsparams}
\centering
\small
\begin{tabular}{lrrr}
\toprule
Variable & $\hat\beta$ & SE & $p$-value \\
\midrule
const & 13.0855 & 0.3860 & 0.0000 \\
org\_type\_corporate & 0.5144 & 0.0679 & 0.0000 \\
org\_type\_research & 0.8841 & 0.2033 & 0.0000 \\
license\_class\_restrictive & -0.2336 & 0.1061 & 0.0276 \\
log\_age & 0.0226 & 0.0334 & 0.4996 \\
log\_params & -0.0252 & 0.0125 & 0.0432 \\
\bottomrule
\multicolumn{4}{l}{$n=1237$ models with parameter counts; $R^2=0.126$.} \\
\end{tabular}
\normalsize
\end{table}


\subsection{Predicting breakout models}
The breakout threshold (90th percentile) is 142,251 downloads.
Table~\ref{tab:logit} gives odds ratios: corporate publishers have $3.89\times$ the odds of breakout versus community (95\% CI $[3.45, 4.39]$, $p<0.001$), research publishers $2.40\times$; restrictive licenses cut the odds roughly in half (OR $0.51$, CI $[0.42, 0.61]$); speech-recognition and image-text-to-text models have the highest task odds (OR $5.77$ and $4.11$). Under repeated stratified 5-fold CV
(10 repeats), the random forest reaches AUC $0.723\pm0.012$ versus $0.716\pm0.012$ for logistic regression, both well above the 0.50 chance level (Table~\ref{tab:cv}); i.e.\ breakout status is modestly but reliably predictable from metadata alone, with the forest only marginally ahead of the linear model. The most important features are
model age dominates ($0.83$ importance), followed by corporate publisher status ($0.06$) (Fig.~\ref{fig:rf}), and the breakout rate is highest among corporate publishers at $27.0\%$ versus $7.6\%$ for community.

\begin{figure}[t]
\centering
\includegraphics[width=\columnwidth]{figures/fig3_rf.png}
\caption{Top: top random-forest features for breakout prediction
(400 trees). Bottom: breakout rate (share of models above the 90th percentile
of downloads) by organization type.}
\label{fig:rf}
\end{figure}

\begin{table}[t]
\caption{Logistic model of breakout adoption (odds ratios; key predictors)}
\label{tab:logit}
\centering
\small
\begin{tabular}{lrrr}
\toprule
Variable & OR & 95\% CI & $p$-value \\
\midrule
const & 0.022 & [0.017, 0.030] & 0.0000 \\
org\_type\_corporate & 3.893 & [3.450, 4.393] & 0.0000 \\
org\_type\_research & 2.403 & [1.871, 3.088] & 0.0000 \\
license\_class\_restrictive & 0.505 & [0.420, 0.608] & 0.0000 \\
tag\_group\_image-text-to-text & 4.113 & [3.424, 4.942] & 0.0000 \\
tag\_group\_automatic-speech-recognition & 5.765 & [4.532, 7.333] & 0.0000 \\
log\_age & 1.132 & [1.082, 1.185] & 0.0000 \\
\bottomrule
\multicolumn{4}{l}{$n=20000$; breakout $\geq$ 142,251 downloads; pseudo-$R^2=0.092$.} \\
\end{tabular}
\normalsize
\end{table}

\begin{table}[t]
\caption{Breakout classifiers: repeated stratified 5-fold CV (10 repeats)}
\label{tab:cv}
\centering
\small
\begin{tabular}{lrr}
\toprule
Model & AUC (mean) & SD \\
\midrule
Random forest (400 trees) & 0.723 & 0.012 \\
Logistic regression & 0.716 & 0.012 \\
Chance (AUC) & 0.500 & n/a \\
Majority-class accuracy & 0.900 & n/a \\
\bottomrule
\end{tabular}
\normalsize
\end{table}


\subsection{Model archetypes}
Silhouette validation selects $k=$ $4$ (silhouette $0.292$) over $k=2,\dots,6$. Table~\ref{tab:archetypes} profiles the clusters: \emph{breakout flagships} (3,445 models; median $\sim$121k downloads, 124 likes), \emph{aging mid-tier} (3,725; $\sim$9.4k downloads, 3.4~years old), \emph{newcomer long tail} (6,239; $\sim$5.0k downloads, under six months old), and \emph{factory long tail} (6,591; publishers averaging $\sim$790 models each but only $\sim$3.8k median downloads (the converter/quantization accounts).

\begin{table}[t]
\caption{K-means model archetypes (median feature values, $k=4$)}
\label{tab:archetypes}
\centering
\small
\begin{tabular}{lrrrrr}
\toprule
Archetype & $n$ & $\log$ dl & Age (yr) & $\log$ org sz & Brk.\% \\
\midrule
Aging mid-tier & 3,725 & 9.15 & 3.44 & 2.08 & 0.0 \\
Newcomer long tail & 6,239 & 8.51 & 0.47 & 1.39 & 0.0 \\
Breakout flagships & 3,445 & 11.70 & 1.17 & 3.99 & 0.0 \\
Factory long tail & 6,591 & 8.24 & 1.51 & 6.67 & 0.0 \\
\bottomrule
\end{tabular}

{\scriptsize Silhouette: $k$=2: 0.26, $k$=3: 0.27, $k$=4: 0.29, $k$=5: 0.28, $k$=6: 0.28.}
\end{table}


\subsection{Centralization and licensing over time}
Fig.~\ref{fig:temporal} tracks cohorts by creation year: the top-ten
organizations' download share falls from 78.6\% for the 2022 cohort to 53.7\% for 2026, while median parameter count rises roughly $100\times$.
License composition (Table~\ref{tab:license}):
permissive licenses dominate at $59.5\%$ overall; the restrictive (behavioral-use) share peaked with the 2023--2024 LLM-license wave ($22.5\%$ and $20.6\%$) before receding to $5.2\%$ in 2026 as permissive licensing rebounded to $67.9\%$.

\begin{table}[h]
\caption{License-class share (\%) by model creation year}
\label{tab:license}
\centering
\small
\begin{tabular}{lrrrrr}
\toprule
Class & 2022 & 2023 & 2024 & 2025 & 2026 \\
\midrule
Permissive & 62.8 & 50.2 & 47.0 & 62.8 & 67.9 \\
Copyleft & 2.4 & 1.7 & 0.5 & 0.6 & 0.5 \\
Restrictive & 4.4 & 22.5 & 20.6 & 7.7 & 5.2 \\
Other & 2.4 & 9.4 & 12.3 & 10.0 & 14.6 \\
Unspecified & 27.9 & 16.2 & 19.6 & 18.8 & 11.8 \\
\bottomrule
\end{tabular}
\normalsize
\end{table}


\begin{figure}[t]
\centering
\includegraphics[width=\columnwidth]{figures/fig4_time.png}
\caption{Top: top-ten organizations' share of cohort downloads by creation
year. Bottom: median parameter count by creation year (detail subsample).}
\label{fig:temporal}
\end{figure}

\section{Discussion}
The central empirical fact is the asymmetry between two levels of
concentration. Across models, downloads are extraordinarily concentrated
(Gini 0.9182); across publishers they are not (HHI 375, below the 1,500
``moderately concentrated'' line). The open-weight ecosystem is thus best
described as a \emph{model-level oligopoly inside an organizationally
competitive market}: thousands of publishers participate, but usage collapses
onto a few artifacts. Several of the top-downloaded models are embedding and
reranking workhorses (not flagship chat models), suggesting that much of
measured ``adoption'' is infrastructural: pipelines pulling the same
checkpoints millions of times rather than end users choosing models.

This qualifies the rebalancing narrative of the recent full-history census
\cite{econ-open-intel}. Community publishers contribute $87.3\%$ of sampled models but only $43.9\%$ of downloads; the typical community model sits deep
in the long tail (median $5{,}856$ downloads) while corporate and research labs
dominate the head. Openness of publishing coexists with concentration of
attention, a pattern familiar from open-source software, now quantified for
model weights.

That breakout status is predictable from metadata alone (random-forest AUC $0.723\pm0.012$ vs.\ 0.50 chance) is both a practical and a cautionary
finding: registries could surface promising models earlier, but the same
features that predict adoption (established org, permissive license, age)
also describe incumbency advantage. The cohort analysis offers modest
optimism: within newer cohorts the top-ten publisher share falls from 78.6\%
(2022) to 53.7\% (2026),
consistent with a slowly diversifying frontier rather than an
entrenching one.

\emph{Limitations.} This is a cross-sectional snapshot (September 29, 2026)
of the top 20,000 models by trailing-30-day downloads: it describes the
head of the distribution at one moment, not the full Hub population of
millions of repos, and rankings reflect current momentum rather than
cumulative adoption. Deleted or gated-then-removed models are invisible
(survivorship bias). Organization type is author-coded from public knowledge,
not API ground truth, with uncoded orgs defaulting to community. Parameter
counts cover only the top-2,000 detail subsample at 61.9\% (many GGUF/quantized repos report no \texttt{safetensors.total}); the full-frame
regressions exclude size. Downloads count pulls, not unique users or
downstream reuse, and the breakout classifiers use contemporaneous features:
they are descriptive, not causal or forecasting claims.

\section{Conclusion}
We measured concentration and adoption in the Hugging Face model ecosystem
from a reproducible September 2026 snapshot of its most-downloaded models.
Three facts stand out. First, the ecosystem is a model-level oligopoly
(Gini $0.918$) inside an organizationally competitive market (HHI $375$):
thousands publish, but usage collapses onto a few artifacts, often embedding
and reranking workhorses rather than flagship chat models. Second, adoption
is structurally legible: publisher type, license class, age, and task explain
$14.8\%$ of log-download variation, breakout status is predictable from
metadata (AUC $0.72$), and, strikingly, parameter count predicts slightly
\emph{fewer} downloads once those factors are held constant, because the
download head is dominated by small, old, infrastructural models. Third, the
frontier diversifies slowly: newer cohorts are less publisher-concentrated
even as models grow roughly $100\times$ larger. For platform designers, the
implication is that discovery, not publishing, is the scarce resource:
$87.3\%$ of models come from community publishers but capture only $43.9\%$
of downloads. For policymakers, ``open weights'' should not be mistaken for
decentralized influence: a handful of checkpoints mediate a large share of
real usage, and the licenses attached to those checkpoints (still $59.5\%$
permissive overall) shape what downstream builders may legally do. Future
work should move from snapshot to panel: linking trailing downloads to
cumulative histories and downstream Space reuse would separate momentum from
installed base.

\begin{thebibliography}{9}
\footnotesize
\bibitem{hf-api} Hugging Face. ``Hugging Face Hub API.''
\url{https://huggingface.co/api/models} (accessed Sept.~29, 2026).
\bibitem{hf-hub-docs} Hugging Face. ``\texttt{huggingface\_hub} documentation:
model info and download statistics.''
\url{https://huggingface.co/docs/huggingface_hub} (accessed Sept.~29, 2026).
\bibitem{econ-open-intel} ``Economies of Open Intelligence: Tracing Power \&
Participation in the Model Ecosystem.'' \textit{Hugging Face Papers},
arXiv:2512.03073, Nov. 2025. [Online]. Available:
\url{https://huggingface.co/papers/2512.03073}
\bibitem{hf-lit-review} ``What do we know about Hugging Face? A systematic
literature review and quantitative validation of qualitative claims.''
arXiv:2406.08205, 2024. [Online]. Available:
\url{https://arxiv.org/pdf/2406.08205v1}
\bibitem{hf-licensing} ``An Empirical Analysis of Machine Learning Model and
Dataset Documentation, Supply Chain, and Licensing Challenges on Hugging
Face.'' arXiv:2502.04484, 2025. [Online]. Available:
\url{https://arxiv.org/pdf/2502.04484}
\bibitem{hf-reuse} ``Beyond Visibility and Technical Reuse: Public Application
Transformation in Open-Source Model Ecosystems.'' arXiv:2607.16687, 2026.
[Online]. Available: \url{https://arxiv.org/pdf/2607.16687}
\bibitem{breiman} L.~Breiman, ``Random forests,'' \textit{Machine Learning},
vol.~45, no.~1, pp.~5--32, 2001.
\bibitem{doj-hhi} U.S.~Dept.~of Justice \& Federal Trade Commission,
``Horizontal Merger Guidelines,'' 2023. (HHI thresholds: unconcentrated below
1500, moderately concentrated 1500--2500, highly concentrated above 2500.)
\end{thebibliography}

\end{document}
